\documentclass[10pt,a4paper]{amsart}
\usepackage[margin=19mm]{geometry}
\usepackage{amsmath,amssymb,mathtools}
\usepackage[scr=rsfs,cal=euler]{mathalfa}
\usepackage{xfrac}
\usepackage[unicode,hidelinks,bookmarks=false]{hyperref}
\newcommand{\ie}{i.e.\ }
\newcommand{\cf}{cf.\ }
\newcommand{\resp}{respectively\ }
\newcommand{\Subsection}{\subsection}
\providecommand{\nobreakdash}{\nobreak\hbox{-}}
\setlength{\emergencystretch}{2em}
\multlinegap0pt
\usepackage{bm}
\usepackage{bbm}
\usepackage{dsfont}
\usepackage{xspace}
\usepackage{cancel}
\usepackage{mathtools}
\usepackage{amsthm}
\makeatletter
\@ifundefined{lemma}{\newtheorem{lemma}{Lemma}}{}
\@ifundefined{proposition}{\newtheorem{proposition}{Proposition}}{}
\@ifundefined{remark}{\newtheorem{remark}{Remark}}{}
\makeatother
\newcommand{\N}{\mathbb{N}}
\newcommand{\Pp}{\mathbb P}
\newcommand{\RR}{\mathbb{R}^2}
\newcommand{\R}{\mathbb{R}}
\newcommand{\Sol}{\operatorname{Sol}}
\newcommand{\TT}{\mathbb{T}^2}
\newcommand{\Z}{\mathbb{Z}}
\newcommand{\bmu}{\boldsymbol{\mu}}
\newcommand{\bomega}{\boldsymbol{\omega}}
\newcommand{\bx}{\mathbf{x}}
\newcommand{\cF}{\mathcal{F}}
\newcommand{\fsol}{Sf}
\newcommand{\hTTf}[1][f]{\mathbf{L}_{#1}}
\newcommand{\hf}{\hat{f}}
\newcommand{\hmu}{\hat{\mu}}
\newcommand{\pcov}{\pi_{\mathrm{cov}}}
\newcommand{\tbeta}{\tilde{\boldsymbol{\eta}}}
\newcommand{\tbmu}{\tilde{\bmu}}
\newcommand{\teta}{\tilde{\eta}}
\newcommand{\tmu}{\tilde{\mu}}
\newcommand{\tx}{\tilde{x}}
\title[Non-uniformly hyperbolic endomorphisms]{Non-uniformly hyperbolic endomorphisms}
\author{Martin Andersson, Pablo D. Carrasco, Radu Saghin}
\date{Source: arXiv:2206.08295v3 (CC BY 4.0)}
\begin{document}
\maketitle

\noindent\emph{Source and licence.} Non-uniformly hyperbolic endomorphisms. Version arXiv:2206.08295v3 (CC BY 4.0), distributed under the Creative Commons Attribution 4.0 International licence, as recorded in the arXiv licence metadata for this exact version and shown on the abstract page https://arxiv.org/abs/2206.08295v3. Full rights notice: licenses/arxiv-2206.08295-CC-BY-4.0.txt.\par\smallskip

\noindent\textit{Editorial scope.} Counted unit: the Remark whose source label is \texttt{None} in the article (source0.tex lines 1500--1502, source label \texttt{None}), delivered together with the article's own complete original proof of it (source0.tex lines 1504--1580, one \texttt{proof} environment of 77 lines). No mathematics is written, repaired or re-derived: statement and proof are the article's own text line for line. Required context retained uncounted: pro:existeexpneg (lines 713--722); le:projective (lines 1341--1349); le:limitproductd (lines 1421--1433); re:product (lines 1471--1473); integrated\_negative\_luap\_exp\_2 (lines 1494--1496); boldmutildextilde (lines 2169--2171); prop:a4 (lines 2197--2199); convergence\_on\_projetive\_bundles (lines 2237--2239). Every definition, display and label of those lines is retained unchanged. Omitted from this package and not redistributed: 1--712, 723--1340, 1350--1420, 1434--1470, 1474--1493, 1497--1499, 1503--1503, 1581--2168, 2172--2196, 2200--2236, 2240--2524. Nothing inside a retained excerpt is omitted or altered: every retained body is delivered exactly as the article's lines are written, so no omission receipt of any kind accompanies this package. Citations: the retained excerpts carry no citation command at all, so no bibliography entry of the article is transcribed and no reference is redistributed. Reference adaptation: none was needed and none was made; every reference inside the delivered unit resolves inside it, so no unresolved reference or citation is produced. Printed slips kept literal: no printed word, symbol, hypothesis or number of the counted statement or of its proof was altered. Layout and class: the article's own class is \texttt{article} with packages including textcomp, amsmath, amsthm, amsfonts, amscd, xy, amssymb, geometry, mathtools, float, subfig, graphicx, caption, babel; the wrapper is the lane's pinned \texttt{amsart} editorial wrapper at 10pt on A4, which restates the theorem-like environments the retained text needs and adds the article's own macro definitions that the retained text needs (\texttt{\textbackslash N}, \texttt{\textbackslash Pp}, \texttt{\textbackslash R}, \texttt{\textbackslash RR}, \texttt{\textbackslash Sol}, \texttt{\textbackslash TT}, \texttt{\textbackslash Z}, \texttt{\textbackslash bmu}, \texttt{\textbackslash bomega}, \texttt{\textbackslash bx}, \texttt{\textbackslash cF}, \texttt{\textbackslash fsol}, \texttt{\textbackslash hTTf}, \texttt{\textbackslash hf}, \texttt{\textbackslash hmu}, \texttt{\textbackslash pcov}, \texttt{\textbackslash tbeta}, \texttt{\textbackslash tbmu}, \texttt{\textbackslash teta}, \texttt{\textbackslash tmu}, \texttt{\textbackslash tx}), with extra wrapper packages (bm, bbm, dsfont, xspace, cancel, mathtools). The delivered numbering is the unit's own. Assets: no retained span contains a figure environment, a graphics-inclusion command, a PDF-page inclusion, a table or a TikZ picture; no image file is copied into this package, the package declares no asset dependency and no image file is redistributed with it. Licence: version arXiv:2206.08295v3 of this article is distributed under the Creative Commons Attribution 4.0 International licence, as recorded in the arXiv licence metadata for this exact version and shown on the abstract page https://arxiv.org/abs/2206.08295v3. A full rights notice is delivered at licenses/arxiv-2206.08295-CC-BY-4.0.txt. Characters: every retained excerpt is pure ASCII, so the delivered engine, XeLaTeX with its default Latin Modern text font, reports no missing character.
\begin{proposition}\label{pro:existeexpneg}
For $\mu$ almost every $x\in\TT$ we have $\chi^{-}_f(x)\leq -C_{\chi}(f)$ and $\chi^{+}_f(x)\geq C_{\chi}(f)+C_{\det}(f)$.
In particular,
\begin{enumerate}
\item
If $C_{\chi}(f)>0$ ($f\in\mathcal U$), then $\chi_f^-(x)<0<\chi_f^+(x)$ for almost every $x\in\TT$;
\item
If $C_{\chi}(f)>-\frac 12C_{\det}(f)$ ($f\in\mathcal U_1$), then $\chi_f^-(x)<\chi_f^+(x)$ for almost every $x\in\TT$.
\end{enumerate}
\end{proposition}

\begin{lemma}[Projectivized cocycles]\label{le:projective}
Assume that $A$ is a two-dimensional cocycle over the invertible map $f:M\rightarrow M$ with invariant measure $\mu$. Let $\Pp f:\Pp M\rightarrow \Pp M$ be its projectivization, which is a bundle map over $f$. Assume that the two Lyapunov exponents $\lambda^+(x)$ and $\lambda^-(x)$ are different at $\mu$-almost every point, with the corresponding Lyapunov subspaces $E^+(x)$ and $E^-(x)$. Then:
\begin{itemize}
\item
If $\mu$ is ergodic then there are exactly two ergodic lifts of $\mu$ to $PM$, $\mu^{p+}$ and $\mu^{p-}$. The disintegrations of $\mu^{p+}$ along the fibers of $\Pp M$ are exactly the Dirac measures at the Lyapunov spaces $E^+$, and a similar statement holds for $\mu^{p-}$.
\item
Suppose that $\mu^p$ is a lift of $\mu$ (not necessarily ergodic) to $\Pp M$, then there is a measurable $f$-invariant function $\rho:M\rightarrow[0,1]$ such that the disintegrations of $\mu^p$ along the fibers of $\Pp M$ are $\rho(x)\delta_{E^+(x)}+(1-\rho(x))\delta_{E^-(x)}$.
\end{itemize}
\end{lemma}

\begin{lemma}[Limits of measures with product disintegrations]\label{le:limitproductd}
Let $X, Y, Z$ be compact metric spaces. Let $\mu^k$ be a sequence of Borel probability measures on $W =X \times Y \times Z$ such that $\mu^k$ converges to $\mu$ in the weak* topology, and for $(\pi_{X})_* \mu^k$-a.e.\@ $x \in X$, the disintegration $\mu_x^k$ is a product measure on $Y \times Z$.\\
Suppose that one of the following two conditions is verified:
\begin{enumerate}

\item
Given any $g \in C(Y, \R)$, the functions $\alpha^k(x) = \int_Y g \ d(\pi_Y)_* \mu_x^k$ and \newline $\alpha(x) =$ $\int_Y g \ d(\pi_Y)_* \mu_x$ can be extended continuously to all of $X$ and, moreover, $\alpha^k(x)$ converges uniformly to $\alpha(x)$. 

\item
The measures $(\pi_{X})_* \mu^k$ are equivalent to $(\pi_{X})_* \mu$ with the Jacobian $J^k=\frac{d(\pi_{X})_* \mu^k}{d(\pi_{X})_* \mu}$ uniformly bounded from above, and $(\pi_Y)_*\mu_x^k$ converges in the weak* topology to $(\pi_Y)_*\mu_x$ for $(\pi_{X})_*\mu$-a.e.\@ $x\in X$.
\end{enumerate}
Then the disintegrations $\mu_x$ of $\mu$ are product measures for $(\pi_{X})_* \mu$-a.e.\@ $x\in X$.
\end{lemma}

\begin{remark}\label{re:product}
The conditions of the Lemma \ref{le:limitproductd} mean in fact that one of the two factors of the product disintegrations (corresponding to M) converges weakly to the factor of the disintegrations of the limit measure (uniformly or pointwise). However the other factor (corresponding to $N$) of the disintegrations may not converge.
\end{remark}

\begin{equation} \label{integrated_negative_luap_exp_2}
\int_{\mathbb T^2}\chi^-(f_n)d\mu=\int_{\Sol}\chi^-(\fsol_n)d\bmu_n= \int_{\Pp \Sol}\|D f_n({\bx}) v\|d\bmu_n^{p-}(\bx,v),
\end{equation}

\begin{equation} \label{boldmutildextilde}
\tilde{\bmu}_{\tx} (C(\omega_1, \ldots \omega_n)) = | \det D (\cF_{\omega_n} \circ \ldots \circ \cF_{\omega_1})(\tx) |
\end{equation}

\begin{proposition}\label{prop:a4}
The measures $\hmu_x^f$, $\bmu_x^f$, $\tilde{\bmu}_{\tx}^f$ $\bmu^f$ and $\tilde{\bmu}^f$ depend continuously on $f$ in the weak* topology in their respective spaces, when seen as maps from $\operatorname{End}_\mu^1(\TT)$. Moreover, the measures $\hmu_x^f$, $\bmu_x^f$ and $\tilde{\bmu}_{\tx}$ vary continuously with $x$ in the weak* topology in their respective spaces, when seen as maps from $\TT$ and $\RR$ respectively.
\end{proposition}

 \begin{proposition}\label{convergence_on_projetive_bundles}
Given a measure $\bmu^{p}$ on $\Pp \Sol$ and a sequence of measures $\bmu_n^{p}$ on $\Pp \Sol$ with corresponding lifts $\tbmu_n^{p}$ and $\tbmu_n^{p}$ on $\Pp (\RR \times \Sigma)$, we have that $\bmu_n^{p}$ converges (weakly*) to $\bmu^{p}$ if and only if $\tbmu_n^{p}$ converges (weakly*) to $\tbmu^{p}$. 
 \end{proposition}

\begin{remark}
We could dispense with $\hTTf$ altogether and replace it with $\Sol$ throughout this paper. Indeed, $\hTTf$ can be seen simply as an embedding of $\Sol$ into $(\TT)^{\Z+}$ with the disadvantages that it depends on $f$ and that its solenoidal structure is somewhat hidden. However, we have decided to stick to $\hTTf$ in all sections but this one, due to its greater familiarity. We hope that this helps the casual reader.
\end{remark}

\begin{proof}[Proof of Theorem B]

Let $(f_n)$ be a sequence in $\mathcal U_1$ converging to $f\in\mathcal U_1$ in the $\mathcal C^1$ topology and let $\bmu_n$ and $\bmu_n^{p-}$ be the lifts to $\Sol$ and $\Pp \Sol$ as above. In virtue of \eqref{integrated_negative_luap_exp_2}, to prove Theorem B it suffices to prove that $\bmu_n^{p-}$ converges weakly* to $\bmu^{p-}$ as $n \to \infty$.
By proposition \ref{convergence_on_projetive_bundles}, this is equivalent to say that $\tbmu_n^{p-}$ converges weakly* to $\tbmu^{p-}$, where $\tbmu_n^{p-}$ and $\tbmu^{p-}$ are the lifts of $\bmu_n^{p-}$ and $\bmu^{p-}$ to $\Pp (\RR \times \Sigma)$, respectively\footnote{Recall that we are denoting $\Sigma=\{1,\cdots, d\}^{\N}$. For the definition of $\Sol$ and the different lifts of invariant measures we send the reader to the \hyperlink{sec:appendix}{Appendix}.}. 

Suppose, for the purpose of contradiction, that $\bmu_n^{p-}$ converges to some $\bmu^p$ different from $\bmu^{p-}$. We are going to show that $\bmu^p$ has product disintegrations, i.e. that it can be written as
\begin{equation} \label{structure_of_bmup} 
\bmu^p = \int_{\TT} \bmu_x \times \nu_x \ d\mu(x) 
\end{equation}
for some (measurable) family of measures $\nu_x$ on $\Pp \RR$. For that we have to analyze the convergence of $\bmu_n^{p-}$ to $\bmu^p$ in some more detail. But instead of analyzing this convergence directly, we are going to analyze the convergence of a sequence of measures on $[0,1]^2 \times \Sigma \times \Pp \RR$ which mirrors this sequence. The advantage of this is that, by doing so, the underlying space is a product space, allowing us to apply Lemma \ref{le:limitproductd}.


Let us write
\[
\begin{array}{lll}
Q = [0,1]^2, & \tbeta_n = \tbmu_n \vert Q \times \Sigma,  & \tbeta_n^{p-} = \tbmu_n^{p-} \vert Q \times \Sigma \times \Pp \RR \\
\teta = \tmu \vert Q, & \tbeta = \tbmu \vert Q \times \Sigma,  & \tbeta^{p} = \tbmu^{p} \vert Q \times \Sigma \times \Pp \RR
\end{array}
\]

Note that the boundary of $Q \times \Sigma \times \Pp \RR$ has zero $\tbmu^{p}$-measure, $(\pi_G,i)_* \tbeta_n^{p-} = \bmu_n^{p-}$, and $(\pi_G,i)_* \tbeta^p = \bmu^p$. (See the Appendix for notation.) We are now in the following situation:

\begin{enumerate}
\item
The projection of $\tbeta_n^{p-}$ to $Q \times \Sigma$ is $\tbeta_n $. The disintegrations of $\tbeta_n^{p-}$ along $ \{(\tx, \bomega) \} \times \Pp \RR$ are $\delta_{E^-_n(\pcov(\tx))}$ (see Lemma \ref{le:projective}). 
\item
The projection of $\tbeta_n$ to $Q$ is $\teta$ (the restriction of Lebesgue to the unit square). The disintegrations of $\tbeta_n$ along $\{\tx\}\times \Sigma$ are $\tbmu_{n,\tx}$ (defined as in \eqref{boldmutildextilde}).
\item
The projection of $\tbeta_n^{p-}$ to $Q$ is also $\teta$. The disintegrations of $\tbeta_n^{p-}$ along $\{\tx\} \times \Sigma \times \Pp \RR$ are $\delta_{E^-_n(\pcov(\tx))}\times\tbmu_{n,\tx}$. (This is because $E^-_n$ is constant on the fibers.)
\item
The disintegrations $\tbmu_{n,\tx}$ vary continuously with respect to  $\tx\in Q$. Furthermore  $\tbmu_{n,\tx}$ converge (uniformly) to $\tbmu_x$ (see Remark \ref{re:product} and Proposition \ref{prop:a4}).
\end{enumerate}

All these considerations show that we are within the hypothesis of Lemma \ref{le:limitproductd}, with $Q$ in place of $X$, $\Sigma$ in place of $Y$, $\Pp \RR$ in place of $Z$, $\tbeta_k^{p-}$ in place of $\mu^k$, and $\tbeta^{p}$ in place of $\mu$. (In fact both condition \emph{1} and condition \emph{2} are satisfied.) As a consequence we have that
\begin{equation} \label{tbtetap_is_intagral_of_products}
\tbeta^p=\int_{ Q}  \tbmu_{\tx}  \times \tilde{\nu}_{\tx} \ d\teta,
\end{equation}
where $\tilde{\nu}_{\tx}$, $\tx \in Q$,  are measures on $\Pp \RR$ (depending measurably on $x$). Writing $\nu_x = \tilde{\nu}_{\tx}$ when $\pcov(\tx) = x$ (which is well-defined $\teta$-a.e.) and applying $(\pi_G,i)_*$ to both sides of \eqref{tbtetap_is_intagral_of_products} gives \eqref{structure_of_bmup}. 

For the remainder of the proof we will abandon the space $\Sol$ in favor of $\hTTf$ and apply Lemma \ref{le:projective}. Thus, writing $\hmu^p = (\phi,i)_* \bmu^p$ and applying $(\phi,i)_*$ to both sides of \eqref{structure_of_bmup} we see that
\[ \hmu^p = \int_{\TT} \hat\mu_x \times \nu_x \ d\mu(x). \]

Since $\bmu_n^p$ are invariant under $PS f_n$, and $PS f_n$ converges to $PS f$, we have that $\bmu^p$ is invariant under $PS f$. Consequently, $\hmu^p$ is invariant under $P \hf$. Lemma \ref{le:projective} tells us that there exists an $\hat f$-invariant function $\hat\rho:\hTTf\rightarrow[0,1]$ such that
$$
\hat\mu^p=\int_{\hTTf}\hat\rho\delta_{E^+}+(1-\hat\rho)\delta_{E^-}\ d\hat\mu.
$$

Since the ergodic components of $\hat f$ are exactly the pre-images under $\pi_{ext}$ of the ergodic components of $f$, we have in fact that $\hat\rho$ is constant on the fibers $\pi_{ext}^{-1}(x)$, in other words there exists a measurable $f$-invariant function $\rho:\mathbb T^2\rightarrow[0,1]$ such that $\hat\rho=\rho\circ\pi_{ext}$. Remembering that $E^-$ is also constant on the fibers, we have that
\begin{eqnarray*}
\int_{\hTTf}\rho\circ\pi_{ext}\delta_{E^+}d\hat\mu&=&\hat\mu^p-\int_{\hTTf}(1-\rho\circ\pi_{ext})\delta_{E^-}d\hat\mu\\
&=&\int_{\mathbb T^2}\nu_x\times\hmu_xd\mu-\int_{\mathbb T^2}(1-\rho)\delta_{E^-}\times\hmu_xd\mu\\
&=&\int_{\mathbb T^2}\left(\nu_x-(1-\rho)\delta_{E^-}\right)\times\hmu_xd\mu\\
&=&\int_{\hTTf}\nu_{\pi_{ext}(\hat x)}-(1-\rho\circ\pi_{ext})\delta_{E^-}d\hat\mu.
\end{eqnarray*}

Then for $\hat\mu$ almost every $\hat x$ we have that 
\[\rho(\pi_{ext}(\hat x))\delta_{E^+(\hat x)}=\nu_{\pi_{ext}(\hat x)}-(1-\rho(\pi_{ext}(\hat x)))\delta_{E^-(\hat x)}.
\]
If $\hat\mu^p$ is different from $\hat\mu^{p-}$ then $\rho$ is nonzero on a set of positive measure. Then there exist a $\mu$-positive measure set of points $x\in\mathbb T^2$ such that $\rho(x)>0$ and $\rho(x)\delta_{E^+(\hat x)}=\nu_{x}-(1-\rho(x))\delta_{E^-(\hat x)}$ holds for $\hmu_x$-a.e $\hat x\in\pi_{ext}^{-1}(x)$. Since the right-hand side depends only on $x$, we have that for $\hmu_x$-a.e. $\hat x$, the unit vector $v^+(\hat x)$ inside $E^+(\hat x)$ is constant, equal to say $v^+(x)$. Since almost every $x\in\TT$ has the property that $\hmu_x$-a.e. point is Lyapunov regular, then there exists a set $A\subset\TT$ of positive $\mu$-measure such that for every $x\in A$ we have
\begin{eqnarray*}
\lim_{n\rightarrow\infty}\frac {I(x,v^+;f^n)}n&=&\lim_{n\rightarrow\infty}\frac 1n\int_{\pi_{ext}^{-1}(x)}\log\|(D_{\hat x}\hat f)^{-n}(v^+)\|d\hmu_x\\
&=&\int_{\pi_{ext}^{-1}(x)}\lim_{n\rightarrow\infty}\frac 1n\log\|(D_{\hat x}\hat f)^{-n}(v^+)\|d\hmu_x\\
&=&\int_{\pi_{ext}^{-1}(x)}-\chi^+d\hmu_x.
\end{eqnarray*}

This implies that $\mu$-a.e. $x\in A$ we have $$\chi^+(x)=-\lim_{n\rightarrow\infty}\frac {I(x,v^+;f^n)}n\leq -C_{\chi}(f).$$

But on the other hand we know from \ref{pro:existeexpneg} that $\mu$-a.e. $x\in\TT$ we have
$$
\chi^-(x)\leq-C_{\chi}(f),
$$
so for $\mu$-a.e. $x\in A$ we have
$$
-2C_{\chi}(f)\geq\chi^+(x)+\chi^-(x)=\lim_{n\rightarrow\infty}\log(\det Df^n(x))\geq C_{\det}(f),
$$
which is a contradiction with the definition of $\mathcal U_1$. This finishes the proof.
\end{proof}

\end{document}
